\begin{thebibliography}{99}
\footnotesize\itemsep=0pt

\bibitem{aefy}
Abe T., Enomoto N., Feigin M., Yoshinaga M., Free reflection multiarrangements
 and quasi-invariants, {i}n preparation.

\bibitem{abe-suy}
Abe T., Suyama D., A basis construction of the extended {C}atalan and {S}hi
 arrangements of the type~{$A_2$}, \href{https://doi.org/10.1016/j.jalgebra.2017.09.024}{\textit{J.~Algebra}} \textbf{493} (2018),
 20--35, \href{https://arxiv.org/abs/1312.5524}{arXiv:1312.5524}.

\bibitem{ath-free}
Athanasiadis C.A., On free deformations of the braid arrangement,
 \href{https://doi.org/10.1006/eujc.1997.0149}{\textit{European~J. Combin.}} \textbf{19} (1998), 7--18.

\bibitem{bm}
Bandlow J., Musiker G., A new characterization for the {$m$}-quasiinvariants of
 {$S_n$} and explicit basis for two row hook shapes, \href{https://doi.org/10.1016/j.jcta.2008.01.011}{\textit{J.~Combin. Theory
 Ser.~A}} \textbf{115} (2008), 1333--1357, \href{https://arxiv.org/abs/0707.3174}{arXiv:0707.3174}.

\bibitem{cha-ves}
Chalykh O.A., Veselov A.P., Commutative rings of partial differential operators
 and {L}ie algebras, \href{https://doi.org/10.1007/BF02125702}{\textit{Comm. Math. Phys.}} \textbf{126} (1990), 597--611.

\bibitem{ede-rei}
Edelman P.H., Reiner V., Free arrangements and rhombic tilings,
 \href{https://doi.org/10.1007/BF02711498}{\textit{Discrete Comput. Geom.}} \textbf{15} (1996), 307--340.

\bibitem{feigin}
Feigin M., Private communication.

\bibitem{fei-ves}
Feigin M., Veselov A.P., Quasi-invariants of {C}oxeter groups and
 {$m$}-harmonic polynomials, \href{https://doi.org/10.1155/S1073792802106064}{\textit{Int. Math. Res. Not.}} \textbf{2002}
 (2002), 521--545, \href{https://arxiv.org/abs/math-ph/0105014}{arXiv:math-ph/0105014}.

\bibitem{fel-ves}
Felder G., Veselov A.P., Action of {C}oxeter groups on {$m$}-harmonic
 polynomials and {K}nizhnik--{Z}amolodchikov equations, \href{https://doi.org/10.17323/1609-4514-2003-3-4-1269-1291}{\textit{Mosc.
 Math.~J.}} \textbf{3} (2003), 1269--1291, \href{https://arxiv.org/abs/math.QA/0108012}{arXiv:math.QA/0108012}.

\bibitem{gpt}
Gao R., Pei D., Terao H., The {S}hi arrangement of the type~{$D_\ell$},
 \href{https://doi.org/10.3792/pjaa.88.41}{\textit{Proc. Japan Acad. Ser.~A Math. Sci.}} \textbf{88} (2012), 41--45,
 \href{https://arxiv.org/abs/1109.1381}{arXiv:1109.1381}.

\bibitem{ot}
Orlik P., Terao H., Arrangements of hyperplanes, \textit{Grundlehren der
 Mathematischen Wissenschaften}, Vol.~300, \href{https://doi.org/10.1007/978-3-662-02772-1}{Springer-Verlag}, Berlin, 1992.

\bibitem{ps-def}
Postnikov A., Stanley R.P., Deformations of {C}oxeter hyperplane arrangements,
 \href{https://doi.org/10.1006/jcta.2000.3106}{\textit{J.~Combin. Theory Ser.~A}} \textbf{91} (2000), 544--597,
 \href{https://arxiv.org/abs/math.CO/9712213}{arXiv:math.CO/9712213}.

\bibitem{sai-log}
Saito K., Theory of logarithmic differential forms and logarithmic vector
 fields, \textit{J.~Fac. Sci. Univ. Tokyo Sect. IA Math.} \textbf{27} (1980),
 265--291.

\bibitem{sai-prim}
Saito K., Period mapping associated to a primitive form, \href{https://doi.org/10.2977/prims/1195182028}{\textit{Publ. Res.
 Inst. Math. Sci.}} \textbf{19} (1983), 1231--1264.

\bibitem{sai-lin}
Saito K., On a linear structure of the quotient variety by a finite reflexion
 group, \href{https://doi.org/10.2977/prims/1195166742}{\textit{Publ. Res. Inst. Math. Sci.}} \textbf{29} (1993), 535--579.

\bibitem{sai-unif}
Saito K., Uniformization of the orbifold of a finite reflection group, in
 Frobenius Manifolds, \textit{Aspects Math.}, Vol.~E36, \href{https://doi.org/10.1007/978-3-322-80236-1_11}{Friedr. Vieweg},
 Wiesbaden, 2004, 265--320.

\bibitem{shi-kl}
Shi J.-Y., The {K}azhdan--{L}usztig cells in certain affine {W}eyl groups,
 \textit{Lecture Notes in Mathematics}, Vol.~1179, \href{https://doi.org/10.1007/BFb0074968}{Springer-Verlag}, Berlin,
 1986.

\bibitem{su-bc}
Suyama D., A basis construction for the {S}hi arrangement of the type
 {$B_\ell$} or {$C_\ell$}, \href{https://doi.org/10.1080/00927872.2013.865051}{\textit{Comm. Algebra}} \textbf{43} (2015),
 1435--1448, \href{https://arxiv.org/abs/1205.6294}{arXiv:1205.6294}.

\bibitem{su-te}
Suyama D., Terao H., The {S}hi arrangements and the {B}ernoulli polynomials,
 \href{https://doi.org/10.1112/blms/bdr118}{\textit{Bull. Lond. Math. Soc.}} \textbf{44} (2012), 563--570,
 \href{https://arxiv.org/abs/1103.3214}{arXiv:1103.3214}.

\bibitem{ter-mult}
Terao H., Multiderivations of {C}oxeter arrangements, \href{https://doi.org/10.1007/s002220100209}{\textit{Invent. Math.}}
 \textbf{148} (2002), 659--674, \href{https://arxiv.org/abs/math.CO/0011247}{arXiv:math.CO/0011247}.

\bibitem{ter-hodge}
Terao H., The {H}odge filtration and the contact-order filtration of
 derivations of {C}oxeter arrangements, \href{https://doi.org/10.1007/s00229-004-0536-z}{\textit{Manuscripta Math.}}
 \textbf{118} (2005), 1--9, \href{https://arxiv.org/abs/math.CO/0205058}{arXiv:math.CO/0205058}.

\bibitem{tsu}
Tsuchida T., On quasiinvariants of {$S_n$} of hook shape, \textit{Osaka~J.
 Math.} \textbf{47} (2010), 461--485, \href{https://arxiv.org/abs/0807.1892}{arXiv:0807.1892}.

\bibitem{yos-prim}
Yoshinaga M., The primitive derivation and freeness of multi-{C}oxeter
 arrangements, \href{https://doi.org/10.3792/pjaa.78.116}{\textit{Proc. Japan Acad. Ser.~A Math. Sci.}} \textbf{78}
 (2002), 116--119, \href{https://arxiv.org/abs/math.CO/0206216}{arXiv:math.CO/0206216}.

\bibitem{yos-char}
Yoshinaga M., Characterization of a free arrangement and conjecture of
 {E}delman and {R}einer, \href{https://doi.org/10.1007/s00222-004-0359-2}{\textit{Invent. Math.}} \textbf{157} (2004), 449--454.

\bibitem{yos-survey}
Yoshinaga M., Freeness of hyperplane arrangements and related topics,
 \href{https://doi.org/10.5802/afst.1413}{\textit{Ann. Fac. Sci. Toulouse Math.}} \textbf{23} (2014), 483--512,
 \href{https://arxiv.org/abs/1212.3523}{arXiv:1212.3523}.

\bibitem{zie}
Ziegler G.M., Multiarrangements of hyperplanes and their freeness, in
 Singularities ({I}owa {C}ity, {IA}, 1986), \textit{Contemp. Math.}, Vol.~90,
 \href{https://doi.org/10.1090/conm/090/1000610}{Amer. Math. Soc.}, Providence, RI, 1989, 345--359.

\end{thebibliography}
